% !TeX root = tcolorbox.tex
% include file of tcolorbox.tex (manual of the LaTeX package tcolorbox)
\clearpage
\section{标签结构}\label{sec:tagging}%
\tcbset{external/prefix=external/tagging_}%

本宏包对标签结构（tagging）提供部分兼容支持。
感谢 Ulrike Fischer，以下内容可以正常工作：

\subsection{已支持的内容}

\begin{itemize}
\item 采用 \refKey{/tcb/capture} 捕获模式 \docValue{minipage} 的基本
  |tcolorbox|（这是标准情形）。\par
  捕获模式 \docValue{hbox} 与 \docValue{fitbox} 暂不支持。
  \refCom{tcbox} 与 \refCom{tcboxfit} 同理。
\item 上部与下部。
\item 文本框的标题。
\item 由 \refCom{newtcolorbox} 等创建的文本框。
\item 定理文本框。根据所用的 \refKey{/tcb/theorem style}，
  标题会被标记为题注（\docValue*{standard}、\docValue*{change standard}、
  \docValue*{plain apart}、\docValue*{change apart}、\docValue*{margin apart}），
  或不被标记（其余样式）。
\item 可断开的 |tcolorbox|。
\item 各类皮肤，如 \refKey{/tcb/enhanced}、\refKey{/tcb/bicolor} 等。
\end{itemize}

\subsection{已知局限}
\begin{itemize}
\item 标签结构代码会在 |tcolorbox| 前后插入额外的 |\par|。
  这可能（将会）与 \refKey{/tcb/before} 和 \refKey{/tcb/after} 的设置相互影响。
  特别是，\refKey{/tcb/nobeforeafter} 将无法按预期工作。
\item 由于标签结构会引入额外的 |\par|，\refEnv{tcbraster} 不能超过一列。
\item 带框标题不会被标记。
\end{itemize}
